\section*{Chapter \ref{ch:ll:logging-capture-and-replay} --- Logging, Capture and Replay}

\begin{solution}{exo:ll:logging-capture-and-replay:1}
$16 + 4 \times 8 = 48$ bytes. The line ``34200000000010 fill 40 of 100 at 999900, position -200'' is 54 characters, 55 with its
newline: about the same. The \omterm{def:ll:logging-capture-and-replay:binlog}{binary log} saves the formatting (the conversions to decimal and the parsing of the format) on
the hot path, not bytes.
\end{solution}

\begin{solution}{exo:ll:logging-capture-and-replay:2}
The decoder reads the arguments with the types recorded for the identifier. Two calls with the same format and different
types (an \texttt{int} and a \texttt{double}) encode different bytes; if they shared an identifier, one of them would be
decoded wrongly.
\end{solution}

\begin{solution}{exo:ll:logging-capture-and-replay:3}
$12 + 189\,875 \times 56 = 10\,633\,012$ bytes; compressed, $2\,831\,089 / 189\,875 \approx 14.9$ bytes per event.
\end{solution}

\begin{solution}{exo:ll:logging-capture-and-replay:4}
$2 \times 10^6 \times 0.02 = 40\,000$ records arrive while the writer is blocked; the ring's capacity is a power of two, so
65\,536 slots, 4 MiB with 64-byte slots.
\end{solution}

\begin{solution}{exo:ll:logging-capture-and-replay:5}
Calibration records written regularly by the writer (not the hot path): pairs of a TSC reading and a reading of a clock
traceable to UTC (disciplined by PTP, chapter~5), close together in time. Offline, each record's TSC value is converted by
interpolating between the calibration pairs around it; the error is that of the reference clock plus the drift between pairs.
\end{solution}

\begin{solution}{exo:ll:logging-capture-and-replay:6}
With the \omterm{def:ll:logging-capture-and-replay:binlog}{binary log}, about $5 \times \qty{5}{\nano\second} = \qty{25}{\nano\second}$ per event, half the event's own work; with
\texttt{snprintf}, about $5 \times \qty{100}{\nano\second} = \qty{500}{\nano\second}$, ten times the event.
\end{solution}

\begin{solution}{exo:ll:logging-capture-and-replay:7}
A statement \texttt{log<\textquotedbl{}order \{\} refused \{\} mask \{\}\textquotedbl{}>(t, cl, code, mask)} in the gate's
refusal path (accepted orders can have their own); replaying \texttt{events.txt} through the gate with a logger, draining and
decoding, the codes counted by reason give the counts on the first line of \texttt{expected.txt}.
\end{solution}

\begin{solution}{exo:ll:logging-capture-and-replay:8}
A decision log says what was decided, not what it was decided from, and cannot be re-run under a debugger or after a fix. The
exchange's recording is not what the process received: it has no gaps, no retransmissions, no local reception times, no
order reports, no parameter changes and no timer firings, and it is in the exchange's order, not the process's. A replay of it
answers what the strategy would have done with perfect data, not what it did.
\end{solution}

\begin{solution}{pb:ll:logging-capture-and-replay:1}
\begin{enumerate}
\item $189\,875 / \qty{120}{\second} \approx 1\,582$ events a second.
\item $1\,582 \times 6.5 \times 3\,600 \approx 37.0$ million events.
\item $37.0 \times 10^6 \times 56 \approx 2.07$ GB raw, about 0.55 GB compressed at 3.76 to 1.
\item The open is among the busiest moments of the day, so the average over it overstates a quiet afternoon; but storage and
replay must be sized for busy days, and event days are busier still.
\item $0.55~\text{GB} \times 252 \times 5 \approx 0.7$ TB.
\item About 139 TB for 200 instruments.
\item The exact binary and its build, its configuration, the parameter snapshots, the reference data (instrument definitions),
the order reports it received, and the replay tools, all versioned with the journal.
\item It holds the outcomes but not their causes: it cannot be re-run, stopped at an event, or used to test a fix.
\item At about 2\,800 times real time, $23\,400 / 2\,800 \approx \qty{8.4}{\second}$: under ten seconds.
\item About half an hour, or less on several cores, since the instruments replay independently.
\item Reading and decoding the journal and the engine's own work per event (memory and the processor); never the market's
timing, since nothing waits.
\item Every decision that depends on time would differ, and the replay would have to wait for real time to pass.
\item \textbf{Named result.} A journal costs 56 bytes per event raw and about 15 compressed (3.76 to 1): at the open's
average of 1\,582 events a second, a 6.5-hour day is 37 million events, 2.1 GB raw and 0.55 GB compressed per instrument, about
0.7 TB over five years; and a replay runs about 2\,800 times faster than real time, a day in under ten seconds.
\item About \qty{5}{\nano\second} with \omterm{def:ll:logging-capture-and-replay:binlog}{deferred formatting}, about \qty{100}{\nano\second} with \texttt{snprintf}.
\item The same actions, bit for bit, up to and including the order in question, with the state that produced it; checked by
comparing the replayed decision log (or its hash) with the production log.
\item It is lost, unless the clock readings are themselves journaled as inputs and replayed.
\item On the engine's thread, where the inputs enter it, in the order it handles them: anywhere earlier the order could
differ, anywhere later something could be missing.
\item That the comparison of two logs finds the first divergence and that nothing before it differs; for a real one, the same
comparison between a production log and its replay, then the journal's event at that point, examined in the replay.
\item One microsecond or better, within 100 microseconds of UTC for high-frequency trading; the TSC's resolution is far finer,
and with calibration against a UTC-traceable clock the log meets both.
\item A complete journal of its inputs, a deterministic program, and the exact build that ran it.
\end{enumerate}
\end{solution}

\begin{solution}{iq:ll:logging-capture-and-replay:1}
It formats text, parses a format, may lock, allocate or write to a file, and its tail is long: a hundred nanoseconds or more
per line. Write a statement identifier, a time stamp and the raw arguments into a ring, let another thread write them out, and
format offline.
\iqlookfor{binary records and \omterm{def:ll:logging-capture-and-replay:binlog}{deferred formatting}.}
\end{solution}

\begin{solution}{iq:ll:logging-capture-and-replay:2}
Identifiers fixed at compile time, a fixed-size record encoded with a few copies into a preallocated single-producer ring per
thread, a writer thread on another core that drains in batches. When the ring is full, drop the record and count it: the hot
path never waits, and the count says what was lost.
\iqlookfor{no work beyond copying, and a drop policy.}
\end{solution}

\begin{solution}{iq:ll:logging-capture-and-replay:3}
Every input as received (market data after the \omterm{def:ll:the-feed-handler:handler}{feed handler}, order reports, timers that depend on time, parameter changes,
control messages) with its reception time, plus the build and configuration; recorded on the engine's thread where the inputs
enter, in its order.
\iqlookfor{inputs as received, time included, at the engine's boundary.}
\end{solution}

\begin{solution}{iq:ll:logging-capture-and-replay:4}
Log the outputs of both with enough detail (actions, or hashes of the state per event), compare them record by record, and go
to the first difference; then look for what differed at that event: time, iteration order, threads, inputs.
\iqlookfor{the first difference, then the cause.}
\end{solution}

\begin{solution}{iq:ll:logging-capture-and-replay:5}
The day's journal, the exact build, its configuration and parameters, and the replay tools; replay to the order, check the
replay against the production log, and read the state that produced the decision. With journals and tools kept, it takes
minutes.
\iqlookfor{journal plus build, and a checked replay.}
\end{solution}

\begin{solution}{iq:ll:logging-capture-and-replay:6}
Each process journals its own inputs in a common binary format with UTC-traceable time stamps, logs its decisions in binary,
ships both to storage compressed; builds and configurations are versioned with them; replay drivers per component run in
continuous integration against recent days and compare with production logs; retention follows the regulation.
\iqlookfor{per-process journals, a common format, versioning, and replay in CI.}
\end{solution}
